\chapter{Giới thiệu về Python}

\section{Python là gì?}

Python là một ngôn ngữ lập trình bậc cao, thông dịch và đa mục đích.
Python được thiết kế với mục tiêu tạo ra một ngôn ngữ có cú pháp rõ ràng,
dễ đọc và dễ học.

Python hỗ trợ nhiều mô hình lập trình, trong đó có lập trình thủ tục,
lập trình hướng đối tượng và lập trình hàm. Nhờ cú pháp đơn giản cùng
hệ sinh thái thư viện phong phú, Python được sử dụng rộng rãi trong
nhiều lĩnh vực như phát triển phần mềm, phát triển Web, khoa học dữ liệu,
trí tuệ nhân tạo, Machine Learning và tự động hóa.

Một chương trình Python đơn giản:

\begin{lstlisting}
print("Hello, Python!")
\end{lstlisting}

Cú pháp của Python có xu hướng sử dụng khoảng trắng để thể hiện cấu trúc
của chương trình. Ví dụ:

\begin{lstlisting}
if age >= 18:
    print("Adult")
else:
    print("Under 18")
\end{lstlisting}

Cách tổ chức này giúp mã nguồn Python có cấu trúc rõ ràng và dễ đọc.

\section{Lịch sử phát triển của Python}

Python được tạo ra bởi Guido van Rossum vào cuối những năm 1980 và
được phát triển trong những năm đầu của thập niên 1990.

Guido van Rossum bắt đầu phát triển Python như một dự án kế thừa từ
ngôn ngữ lập trình ABC. Mục tiêu của ông là xây dựng một ngôn ngữ có
cú pháp đơn giản, dễ sử dụng nhưng vẫn đủ mạnh để phát triển các
chương trình thực tế.

Phiên bản Python đầu tiên được công bố rộng rãi là Python 0.9.0,
phát hành vào năm 1991. Phiên bản này đã có nhiều đặc điểm quan trọng
của Python hiện đại, bao gồm các kiểu dữ liệu như danh sách (list),
chuỗi (string), từ điển (dictionary), hàm và cơ chế xử lý ngoại lệ.

\subsection{Python 1.x}

Python 1.0 được phát hành vào năm 1994. Trong giai đoạn này, ngôn ngữ
tiếp tục được phát triển và bổ sung nhiều tính năng quan trọng.

\subsection{Python 2.x}

Python 2.0 được phát hành vào năm 2000. Đây là một cột mốc quan trọng
trong lịch sử Python, với nhiều cải tiến về ngôn ngữ và khả năng xử lý
dữ liệu.

Python 2 tiếp tục được sử dụng trong nhiều năm. Tuy nhiên, Python 2
cuối cùng đã được ngừng hỗ trợ chính thức vào năm 2020.

\subsection{Python 3.x}

Python 3.0 được phát hành vào năm 2008. Đây là một phiên bản được
thiết kế để cải thiện ngôn ngữ và khắc phục một số hạn chế của Python 2.

Python 3 không hoàn toàn tương thích ngược với Python 2. Vì vậy, quá
trình chuyển đổi từ Python 2 sang Python 3 diễn ra trong nhiều năm.

Ngày nay, Python 3 là dòng phiên bản chính được sử dụng và phát triển.

\section{Guido van Rossum -- tác giả của Python}

Guido van Rossum là lập trình viên người Hà Lan và là người tạo ra
ngôn ngữ lập trình Python.

Ông bắt đầu phát triển Python vào cuối năm 1989 tại Centrum Wiskunde
\& Informatica (CWI) ở Hà Lan. Phiên bản đầu tiên của Python được công
bố vào năm 1991.

Guido van Rossum giữ vai trò chính trong việc định hướng sự phát triển
của Python trong nhiều năm. Ông từng được cộng đồng Python gọi bằng
danh hiệu ``Benevolent Dictator For Life'', thường được viết tắt là
BDFL.

Năm 2018, Guido van Rossum tuyên bố từ bỏ vai trò BDFL. Sau đó, quá
trình định hướng phát triển Python được chuyển sang mô hình quản trị
dựa trên Python Steering Council.

\section{Tên gọi Python}

Tên ``Python'' không xuất phát từ loài rắn như nhiều người thường nghĩ.
Guido van Rossum đặt tên ngôn ngữ theo chương trình hài kịch truyền hình
của Anh có tên \textit{Monty Python's Flying Circus}.

Tên gọi ngắn gọn ``Python'' được lựa chọn vì Guido van Rossum là người
yêu thích chương trình này.

\section{Python ngày nay}

Python hiện được sử dụng trong rất nhiều lĩnh vực của công nghệ thông
tin. Một số lĩnh vực tiêu biểu bao gồm:

\begin{itemize}
    \item Phát triển phần mềm.
    \item Phát triển Web.
    \item Khoa học dữ liệu.
    \item Trí tuệ nhân tạo.
    \item Machine Learning.
    \item Deep Learning.
    \item Tự động hóa.
    \item Tính toán khoa học.
\end{itemize}

Đặc biệt, hệ sinh thái Python có nhiều thư viện mạnh như NumPy, Pandas,
Matplotlib, Scikit-learn, TensorFlow và PyTorch. Đây là một trong những
yếu tố quan trọng giúp Python trở thành ngôn ngữ phổ biến trong khoa học
dữ liệu và Machine Learning.

\section{Tóm tắt chương}

Trong chương này, chúng ta đã tìm hiểu:

\begin{itemize}
    \item Khái niệm cơ bản về Python.
    \item Lịch sử hình thành và phát triển của Python.
    \item Các giai đoạn Python 1.x, Python 2.x và Python 3.x.
    \item Guido van Rossum -- người tạo ra Python.
    \item Nguồn gốc tên gọi Python.
    \item Một số lĩnh vực sử dụng Python hiện nay.
\end{itemize}
